\section{极限计算与代数化简}
\label{sec:02-Calculus/02-Limits-and-Continuity/03}

在计算机里运行双精度浮点数计算 \(f(x)=\frac{\sqrt{1+x}-1}{x}\) 时，会遇到一个诡异的现象：数学上当 \(x\to0\) 时该式趋于 \(0.5\)；但实测发现，\(x=10^{-8}\) 时结果还算正常，\(x=10^{-12}\) 时开始偏离，到了 \(x=10^{-17}\) 时直接返回了 \(0\)。

只要将表达式乘上共轭因子进行改写：

\[
\frac{\sqrt{1+x}-1}{x}=\frac{1}{\sqrt{1+x}+1},\qquad x\neq0.
\]

在相同的机器上跑右侧的表达式，\(x=10^{-17}\) 时依然能精确给出 \(0.5\)。两者的根本区别在于：左式在小参数下需要将两个极其接近的数相减，有效数字在这一步被彻底吃光（灾难性抵消）；而右式从结构上消除了相减动作。

代数化简因此具备双重意义：在\textbf{数值计算}中，它是规避舍入误差与有效数字丧失的工程手段；在\textbf{极限理论}中，它是解除 \(\frac{0}{0}\) 型未定式的核心技巧。这两者依赖的是同一个恒等变形，成立的前提条件也完全一致：\(x\neq0\)。

本节的重点在于探讨：哪些改写与拆分是合法有效的？背后的逻辑是什么？以及当代数变形彻底失效时，数学还留下了什么工具？

\subsection{四则运算法则：避免回归原始定义}

极限的 \(\varepsilon\)-\(\delta\) 定义是事后验证的工具，若计算每个极限都从定义出发，效率极低。好在极限的承诺结构具有良好的可组合性：若已知 \(f(x)\to L\) 且 \(g(x)\to M\)（其中 \(L,M\) 均为有限实数），则：

\[
 f+g\to L+M,\qquad fg\to LM,\qquad cf\to cL,
\]

且当 \(M\neq0\) 时，有 \(f/g\to
 L/M\)。

这些运算法则的证明逻辑完全一致：面对给定的 \(\varepsilon\)，将误差指标拆分并分派给各个子项——加法分配 \(\varepsilon/2\)，乘法按各自的界分配——最后取各自 \(\delta\) 的交集（即最小值）。极限可以像代数一样进行四则运算，本质上是因为这种``责任拆分与外包''在逻辑上永远成立。

这也清晰地揭示了四则法则的使用前提：\textbf{每个子项必须各自收敛到一个有限的极限值}。一旦出现 \(\infty-\infty\) 或 \(0\cdot\infty\) 这类形式，由于子项本身没有确定的有限目标值可以承接责任，拆分法则便彻底失效。

因此，多项式可以直接代入求极限；有理函数在分母极限非零处也可以直接代入。需要注意的是，``直接代入''并非一条独立的数学法则，它是极限四则法则与函数连续性结合后的自然推论。

\subsection{未定式：信息缺失的结构标记}

当代入极限得到 \(\frac{0}{0}\) 或 \(\frac{\infty}{\infty}\) 时，这类形式被称为\textbf{未定式}（Indeterminate Form）。``未定''并不意味着``答案不存在''或``随机不可测''，而是指\textbf{仅凭``分子分母同时趋于零''这一粗粒度信息，不足以判断整个商的极限行为}。

观察以下三个同样呈现 \(\frac{0}{0}\) 形式的例子在 \(x\to0\) 时的表现：

\[
\frac{x}{x}=1,\qquad \frac{x^2}{x}=x\to0,\qquad \frac{x}{x^2}=\frac{1}{x}\to\infty.
\]

三者分别表现为恒定常数、趋于零以及无界发散。结构完全不同的函数可以套着同一个未定式的壳子，因此必须拆解内部的微观结构。

拆解的具体手段因题而异，但核心目标完全一致：\textbf{在去除目标点自身后的局部邻域（即穿孔邻域 \(0<\lvert x-a\rvert<\delta\)）内，将表达式改写为非未定形式}。例如 \(\frac{x^2-1}{x-1}\) 依靠约去公因子，根式表达依靠乘共轭因子。

这种变形的合法性源于极限的基本定义：\textbf{极限只审视目标点附近的趋势，完全不在乎目标点自身的取值}。只要改写后的表达式在 \(x\neq a\) 时与原式严格相等，该变形在极限运算中就是完全合法的。

\subsection{夹逼定理：将分析责任转包出去}

代数变形并非万能。以 \(\frac{\sin x}{x}\) 在 \(x\to0\) 时的极限为例，这同样是一个 \(\frac{0}{0}\) 型未定式，但正弦函数既无法进行多项式因式分解，也无法通过共轭根式化简。

此时需要转换思路：若无法直接计算该函数，可以通过寻找上界与下界将其``押送''至目标值。

\begin{theorem}[夹逼定理 / Squeeze Theorem]
若在点 \(a\) 的某个穿孔邻域内成立 \(g(x)\le f(x)\le h(x)\)，且 \(\lim_{x\to a}g(x)=L\) 以及 \(\lim_{x\to a}h(x)=L\)，则 \(\lim_{x\to a}f(x)=L\)。
\end{theorem}

从承诺机制的角度来看，夹逼定理实现了``责任转移''。面对给定的 \(\varepsilon\)，你无需直接分析复杂的 \(f(x)\)，只需将 \(\varepsilon\) 转发给上下界函数 \(g(x)\) 和 \(h(x)\)，分别获取对应的 \(\delta_1\) 与 \(\delta_2\)。取 \(\delta=\min(\delta_1,\delta_2)\)，在该范围内，\(g(x)\) 与 \(h(x)\) 均被约束在 \(L\pm\varepsilon\) 的区间内，夹在中间的 \(f(x)\) 必然被锁定在相同范围内。

夹逼定理的应用有两个关键约束条件：

1. \textbf{两侧边界必须收敛至同一个极限值}。若 \(g\to0\) 而 \(h\to2\)，中间的 \(f(x)\) 仍有自由游荡的空间。
2. \textbf{必须具备双侧约束}。仅有单侧界只能说明函数被压在某条曲线下方，无法锁定极限。

夹逼定理最典型的应用场景，是一个因子趋于零、另一个因子有界但无极限的情形。例如 \(x\sin(1/x)\) 在 \(x\to0\) 时的极限：\(\sin(1/x)\) 本身在 \(\pm1\) 之间无限振荡、极限不存在，但满足：

\[
-\lvert x\rvert\le x\sin\frac{1}{x}\le \lvert x\rvert.
\]

由于两侧边界在 \(x\to0\) 时均趋于零，夹在中间的函数极限必然为零。

\begin{center}
\begin{tikzpicture}[x=3.6cm,y=3.6cm,font=\small,>=stealth]
\draw[->] (-1.15,0) -- (1.2,0) node[right] {\(x\)};
\draw[->] (0,-1.15) -- (0,1.2) node[above] {\(y\)};
\draw[red!65!black,dashed,thick] (-1.05,1.05) -- (1.05,-1.05);
\draw[red!65!black,dashed,thick] (-1.05,-1.05) -- (1.05,1.05);
\draw[blue!70!black,thick] plot[domain=0.07:1.05,samples=300] (\x,{\x*sin(deg(1/\x))});
\draw[blue!70!black,thick] plot[domain=-1.05:-0.07,samples=300] (\x,{\x*sin(deg(1/\x))});
\node[right,red!65!black,font=\footnotesize] at (1.02,1.06) {\(y=\lvert x\rvert\)};
\node[right,red!65!black,font=\footnotesize] at (1.02,-1.06) {\(y=-\lvert x\rvert\)};
\end{tikzpicture}

\emph{两条虚线构成一个在原点收拢的漏斗。中间的曲线在 \(x\to0\) 时震荡频率无限增加，但其振幅被边界严格压制，极限收敛于零。}
\end{center}

未乘 \(x\) 之前，\(\sin(1/x)\) 因无限振荡导致极限不存在；乘以 \(x\) 之后，振荡虽然依旧存在，但振幅受到包络线的约束而衰减。极限的存在性取决于振幅是否受到有效压制。

---

\subsection{重要极限 \texorpdfstring{\(\sin x/x\)}{sin x / x} 的几何推导}

回到 \(\frac{\sin x}{x}\) 的求解。在单位圆中比较三角形面积、扇形面积与外切三角形面积，对于 \(x\in(0,\pi/2)\)，可建立如下不等式链：

\[
\cos x\le \frac{\sin x}{x}\le1.
\]

当 \(x\to0^+\) 时，左端 \(\cos x\to1\)，右端为常数 \(1\)，根据夹逼定理，中间项的右极限为 \(1\)。结合函数的偶函数对称性，可得：

\[
\lim_{x\to0}\frac{\sin x}{x}=1.
\]

该极限是微积分的核心基石之一，正弦与余弦函数的求导公式均以此为基础。

需要特别指出的是：\textbf{上述推导严格依赖于弧度制}。弧度制保证了单位圆弧长与扇形面积公式的简洁性。若采用角度制，不等式中的系数将发生改变，极限值将变为 \(\pi/180\)。单位体系是数学推导的基础组成部分，在数值代码中混用角度与弧度会导致严重的计算错误。

---

\subsection{无穷远处的极限：主导尺度提取}

另一类常见情形是自变量趋于无穷。例如求 \(\lim_{x\to\infty}\frac{2x^2-x+1}{5x^2+3}\)，此时分子与分母均发散，呈现 \(\frac{\infty}{\infty}\) 形式，直接拆分法则失效。

此时的关键在于提取主导尺度：分子分母同除以最高次项 \(x^2\)，将式子改写为：

\[
\frac{2-1/x+1/x^2}{5+3/x^2}.
\]

改写后，每一项在 \(x\to\infty\) 时均有确定的极限：分子趋于 \(2\)，分母趋于 \(5\)，整体极限为 \(2/5\)。

对于有理函数（多项式之比），可总结出一般规律：

\begin{itemize}
\tightlist
\item 分子次数低于分母，极限为 \(0\)；
\item 分子分母次数相同，极限为最高次项系数之比；
\item 分子次数高于分母，极限无界（发散）。
\end{itemize}

比起记忆上述结论，更重要的是掌握``同除以主导尺度''的核心动作。在包含根式或指数项的复杂表达式中，直接忽略低次项容易出错，唯有通过同除主导项、将未定式转化为收敛项，才能确保推导的严密性。

\subsection{总结与边界}

本节提供了处理极限的三种基本手段：

\begin{enumerate}
\tightlist
\item \textbf{四则运算法则}：适用于各子项均收敛至有限值的情形；
\item \textbf{恒等变形}：适用于可通过代数化简解除未定式的情形；
\item \textbf{夹逼定理}：适用于目标函数难以直接分析、但可构造确切上下界的情形。
\end{enumerate}

当上述方法均失效时（例如 \(\lim_{x\to0}\frac{e^x-1-x}{x^2}\) 这类需要更高阶局部展开的式子），则需要借助泰勒展开（Taylor Series）与洛必达法则（L'Hôpital's Rule）等更高级的工具。

最后需要明确数学理论与工程实现的边界：\textbf{数学函数的极限性质不等同于浮点数计算的实际行为}。数学上完全等价的表达，在计算机有限精度浮点数表示下可能产生截然不同的数值结果。确定数学极限的存在性是分析学问题，而保障数值代码在极限邻域的稳定性则是数值分析的范畴。

在目前的讨论中，极限严格排除了目标点 \(a\) 本身。下一节将解除这一限制：探究当函数在目标点有定义、且该点函数值恰好等于周围趋势的极限值时，系统将具备哪些优良的数学性质。